% tuto_mpm_beton.tex -- Tutoriel : reproduire le comportement du béton armé en MPM
% Compiler avec pdfLaTeX (Overleaf : compilateur par défaut). Les images tuto_beton_*.png doivent
% se trouver dans le même dossier.
\documentclass[11pt,a4paper]{article}

\usepackage[utf8]{inputenc}
\usepackage[T1]{fontenc}
\usepackage{lmodern}
\usepackage[french]{babel}
\usepackage[margin=2.2cm]{geometry}
\usepackage{amsmath,amssymb}
\usepackage{xcolor}
\usepackage{booktabs}
\usepackage{graphicx}
\usepackage{tikz}
\usetikzlibrary{arrows.meta,positioning,calc}
\usepackage{listings}
\usepackage{hyperref}
\hypersetup{colorlinks=true,linkcolor=blue!60!black,urlcolor=blue!60!black}

% ---------------------------------------------------------------- style des listings
\definecolor{codebg}{RGB}{247,247,250}
\definecolor{codekw}{RGB}{0,80,160}
\definecolor{codecomment}{RGB}{90,120,90}
\definecolor{codestring}{RGB}{160,60,20}
\lstset{
  language=Python,
  basicstyle=\ttfamily\small,
  keywordstyle=\color{codekw}\bfseries,
  commentstyle=\color{codecomment}\itshape,
  stringstyle=\color{codestring},
  backgroundcolor=\color{codebg},
  frame=single, framerule=0pt, rulecolor=\color{codebg},
  xleftmargin=6pt, xrightmargin=6pt,
  breaklines=true, breakatwhitespace=true,
  showstringspaces=false, tabsize=4,
  columns=fullflexible, keepspaces=true,
  literate=
    {é}{{\'e}}1 {è}{{\`e}}1 {ê}{{\^e}}1 {ë}{{\"e}}1
    {à}{{\`a}}1 {â}{{\^a}}1 {î}{{\^i}}1 {ï}{{\"i}}1
    {ô}{{\^o}}1 {ù}{{\`u}}1 {û}{{\^u}}1 {ç}{{\c{c}}}1
    {É}{{\'E}}1 {È}{{\`E}}1 {À}{{\`A}}1 {œ}{{\oe}}1
}
\newcommand{\code}[1]{\texttt{\detokenize{#1}}}
\newcommand{\bm}[1]{\boldsymbol{#1}}
\newcommand{\tr}{\operatorname{tr}}

\title{Le béton armé en MPM :\\ reproduire le comportement d'une poutre de la Rance}
\author{}
\date{}

\begin{document}
\maketitle

\begin{abstract}
\noindent Un solide MPM dont on se contente de régler un seuil d'endommagement à 5~\% et une rupture
à 20~\% de déformation se comporte comme un caoutchouc : il fléchit de façon spectaculaire, puis
se pulvérise. Le béton, lui, fissure vers $10^{-4}$, s'écrase vers $4\cdot10^{-3}$, résiste
treize fois moins bien en traction qu'en compression, et ne travaille jamais seul : il est armé. Ce
tutoriel reprend les lois de comportement simplifiées de l'article de Cremona et Houde sur les
poutres de la Rance (loi de Desayi en compression, raidissement en traction, acier bilinéaire,
résistances aléatoires) et les traduit en MPM : endommagement unilatéral par direction principale,
armatures noyées dans les particules, essai de flexion 4 points piloté en déplacement. Le résultat
est comparé à une analyse de section à la manière de l'article. On y décrit aussi ce qui n'a pas
marché.
\end{abstract}

\tableofcontents
\bigskip

% ================================================================
\section{Pourquoi le solide paraissait flasque}

Le solveur du tutoriel précédent (\code{tuto_mpm_solide}) utilise une loi corotationnelle avec un
seuil d'endommagement $\varepsilon_0 = 0.05$ et une rupture à $\varepsilon_f = 0.20$, un module
d'Young sans unité, et on le charge en multipliant la gravité. Comparé à l'article, c'est un autre
matériau :

\begin{center}
\begin{tabular}{@{}lll@{}}
\toprule
 & modèle précédent & béton (article, poutre 421) \\
\midrule
début de l'endommagement en traction & $0.05$ & $\varepsilon_{cr} = f_{cr}/E_c \approx 1.2\cdot10^{-4}$ \\
rapport au modèle précédent & & $\div 430$ \\
déformation d'écrasement & $0.20$ (rupture) & $\varepsilon_{cu} = 4\cdot10^{-3}$ ($\div 50$) \\
résistance traction / compression & $1$ (critère symétrique) & $3.5/47.1 \approx 0.074$ \\
après fissuration & la raideur tend vers 0 & l'acier prend le relais (plateau) \\
chargement & $g$ multiplié par un facteur & vérins en déplacement imposé \\
\bottomrule
\end{tabular}
\end{center}

Un béton qui attendrait $5$~\% de déformation avant de s'endommager n'existe pas : c'est ce qui
donnait la sensation de \og flasque \fg{}. Il ne suffit donc pas de raidir le module : il faut des
seuils réalistes, une loi propre à chaque sens de sollicitation, des armatures, et un essai
piloté comme en laboratoire.

% ================================================================
\section{Ce que l'article apporte}

L'article modélise des poutres en béton armé ou précontraint (section de $200\times200$~mm) avec
des outils de résistance des matériaux : analyse de section (couches de béton et d'acier), puis
intégration des courbures le long de la poutre. On en retient les ingrédients suivants.

\paragraph{Béton en compression (Desayi/Neville).} Avec $\varepsilon_{ult} = 1.8\,f_c/E_c$ la
déformation au pic,
\begin{equation}
  \sigma = \frac{-2\,\varepsilon\, f_c}{\varepsilon_{ult}\,\big(1 + (\varepsilon/\varepsilon_{ult})^2\big)},
  \qquad \varepsilon_{cu} = -0.004 .
  \label{eq:desayi}
\end{equation}

\paragraph{Béton en traction (raidissement).} Linéaire jusqu'à la résistance $f_{cr}$, puis
\begin{equation}
  \sigma = \frac{f_{cr}}{1 + \sqrt{500\,\varepsilon}} .
  \label{eq:tension}
\end{equation}
Cette décroissance lente représente l'apport du béton qui entoure les barres entre les
fissures (\emph{tension stiffening}, Collins et Mitchell).

\paragraph{Acier.} Loi bilinéaire, élastique jusqu'à $f_y$, puis parfaitement plastique, symétrique
en traction et en compression, jusqu'à $\pm0.02$ (au-delà, l'article considère la barre rompue).

\paragraph{Aléa.} Les résistances sont tirées au hasard (Monte-Carlo, lois lognormales
indépendantes pour rester positives). Le tableau~3 de l'article donne moyenne et écart type de
chaque poutre.

\paragraph{Essai.} Flexion 4 points sur appuis simples : portée $L = 2.00$~m, charges espacées de
$x_p = 0.60$~m des appuis (moment constant sur $0.80$~m au centre). Les poutres du type 4 ont
quatre lits d'armatures (tableau~2 de l'article) ; on prend ici les données de la poutre 421
($f_c = 47.1$~MPa, $f_{cr} = 3.5$~MPa, $E_c = 30.4$~GPa).

\paragraph{Ce que l'article ne fait pas} (et nous non plus) : les cycles de chargement sont
traités comme un chargement monotone, la corrosion est laissée de côté ici, et la précontrainte n'est
pas reprise (poutre passive).

% ================================================================
\section{De la résistance des matériaux au MPM}

La RDM raisonne sur une \emph{section} et une loi uniaxiale $\sigma(\varepsilon)$ ; le MPM travaille
sur des \emph{particules} qui portent un gradient de déformation $\mathbf F$ (2D). Il faut donc
traduire chaque ingrédient :

\begin{center}
\begin{tabular}{@{}p{5.3cm}p{9.5cm}@{}}
\toprule
Article (RDM) & Ici (MPM) \\
\midrule
couches de béton d'une section & particules ; les lois s'appliquent par \emph{direction principale} \\
loi non linéaire $\sigma(\varepsilon)$ & raideur \emph{sécante} $r(\kappa)$ pilotée par la plus grande déformation $\kappa$ atteinte (historique, irréversible) \\
traction et compression différentes & deux facteurs $r_t$ et $r_c$, choisis selon le signe de la contrainte principale \\
barres d'armature & fraction volumique d'acier $\phi$ dans une rangée de particules, loi de barre uniaxiale \\
effet raidisseur & on garde la loi~\eqref{eq:tension} telle quelle \\
Monte-Carlo & résistances lognormales tirées \emph{par particule} \\
essai 4 points & appuis glissants, vérins en vitesse imposée, mesure de la réaction $P$ \\
\bottomrule
\end{tabular}
\end{center}

% ================================================================
\section{Le modèle de matériau}

\subsection{Contraintes élastiques principales}

À chaque pas, on décompose $\mathbf F = \mathbf U\,\mathrm{diag}(s_1, s_2)\,\mathbf V^{\mathsf T}$ (SVD,
\code{ti.svd}) et on prend les déformations de Hencky $e_i = \ln s_i$. La contrainte élastique
principale (déformations planes, coefficients de Lamé $\mu$ et $\lambda$) est
\begin{equation}
  \sigma_i^{\rm el} = 2\mu\, e_i + \lambda\,(e_1 + e_2),\qquad i = 1, 2 ,
\end{equation}
et la contrainte de Kirchhoff déposée sur la grille vaut $\bm\tau = \sum_i \sigma_i\,\mathbf u_i\mathbf u_i^{\mathsf T}$
où $\mathbf u_i$ sont les colonnes de $\mathbf U$. Sans endommagement, $\sigma_i = \sigma_i^{\rm el}$.

\subsection{Déformation équivalente : un critère en contrainte}

Pour décider si le béton est sollicité en traction ou en compression, on utilise
\begin{equation}
  \varepsilon_t = \frac{\max(\sigma_1^{\rm el},\, \sigma_2^{\rm el},\, 0)}{E_c},
  \qquad
  \varepsilon_c = \frac{\max(-\sigma_1^{\rm el},\, -\sigma_2^{\rm el},\, 0)}{E_c}.
\end{equation}
C'est un critère de Rankine écrit en déformation : sous compression uniaxiale, la contrainte
transversale est nulle, donc $\varepsilon_t = 0$. Un critère sur la plus grande \emph{déformation}
principale détecterait de la traction dans une zone comprimée, par effet Poisson (le béton comprimé
s'allonge dans la direction transverse), et risquerait de produire de fausses fissures sous les
vérins (choix de conception, non comparé). On ajoute l'historique irréversible
\begin{equation}
  \kappa_t \leftarrow \max(\kappa_t, \bar\varepsilon_t),\qquad
  \kappa_c \leftarrow \max(\kappa_c, \bar\varepsilon_c),
\end{equation}
où la barre désigne la valeur lissée sur la grille (\S\ref{sec:nonlocal}).

\subsection{Raideurs sécantes}

Les lois~\eqref{eq:desayi} et~\eqref{eq:tension} deviennent des facteurs multiplicatifs
$r = \sigma_{\rm loi}(\kappa)/(E_c\,\kappa)$ de la contrainte élastique. En traction, avec
$\varepsilon_{cr} = f_{cr}/E_c$ :
\begin{equation}
  r_t(\kappa) =
  \begin{cases}
    1 & \kappa \le \varepsilon_{cr},\\[2pt]
    \min\!\Big(1,\ \dfrac{f_{cr}}{(1+\sqrt{500\,\kappa})\,E_c\,\kappa}\Big) & \kappa > \varepsilon_{cr},
  \end{cases}
  \qquad r_t \ge r_{\min} = 10^{-3}.
\end{equation}
En compression, avec $x = \kappa/\varepsilon_{ult}$ :
\begin{equation}
  r_c(\kappa) = \min\!\Big(1,\ \frac{2/1.8}{1 + x^2}\Big),
\end{equation}
puis, au-delà de $\varepsilon_{cu} = 4\cdot10^{-3}$, une décroissance linéaire jusqu'à
$r_{\rm crush} = 0.1$ sur $[\varepsilon_{cu},\,1.5\,\varepsilon_{cu}]$. Deux remarques :
\begin{itemize}
  \item la tangente initiale de la loi de Desayi vaut $2f_c/\varepsilon_{ult} \approx 1.11\,E_c$,
        donc supérieure à $E_c$ : le \code{min(1, ...)} conserve un début linéaire de module exact $E_c$ ;
  \item l'article coupe la contrainte à $\varepsilon_{cu}$ ; on adoucit sur $[\varepsilon_{cu}, 1.5\varepsilon_{cu}]$
        pour ne pas créer de saut de contrainte d'une particule à l'autre.
\end{itemize}
Valeurs (poutre 421, $E_c = 30.4$~GPa, $f_c = 47.1$ et $f_{cr} = 3.5$~MPa) :

\begin{center}
\begin{tabular}{@{}rrr@{\qquad}rrr@{}}
\toprule
\multicolumn{3}{c}{traction} & \multicolumn{3}{c}{compression} \\
$\kappa$ & $\sigma$ (MPa) & $r_t$ & $\kappa$ & $\sigma$ (MPa) & $r_c$ \\
\midrule
$2\cdot10^{-4}$ & 2.66 & 0.437 & $10^{-3}$ & 29.9 & 0.985 \\
$5\cdot10^{-4}$ & 2.33 & 0.154 & $2\cdot10^{-3}$ & 44.6 & 0.734 \\
$10^{-3}$ & 2.05 & 0.067 & $2.79\cdot10^{-3}$ (pic) & 47.1 & 0.556 \\
$1.5\cdot10^{-3}$ & 1.88 & 0.041 & $4\cdot10^{-3}$ & 44.2 & 0.363 \\
$10^{-2}$ & 1.08 & 0.004 & $5\cdot10^{-3}$ & 27.6 & 0.182 \\
\bottomrule
\end{tabular}
\end{center}

\subsection{Unilatéralité : les fissures se referment}

Chaque contrainte principale reçoit le facteur de \emph{son} signe :
\begin{equation}
  \sigma_i = \begin{cases} r_t(\kappa_t)\ \sigma_i^{\rm el} & \sigma_i^{\rm el} > 0,\\[2pt]
                            r_c(\kappa_c)\ \sigma_i^{\rm el} & \sigma_i^{\rm el} \le 0. \end{cases}
\end{equation}
Une particule fissurée en traction ($r_t \approx 0$) garde donc sa raideur en compression : la
fissure se referme quand on recharge dans l'autre sens, comme dans le béton réel.

\subsection{Lissage non local}
\label{sec:nonlocal}

On reprend sans changement l'idée du tutoriel précédent : $\varepsilon_t$ et $\varepsilon_c$ sont
déposés sur la grille (moyenne pondérée par les poids B-spline) puis relus, ce qui donne une
déformation lissée sur environ cinq cellules. Sans ce lissage, chaque particule s'endommage seule et
la fissure n'a pas de largeur. Ici la largeur vaut quelques $\Delta x = 1.56$~cm, du même ordre que
l'espacement des fissures de flexion sur une section de $20$~cm.

\subsection{Les armatures}

Les barres d'un lit ne sont pas maillées : la rangée de particules située à leur hauteur porte une
fraction volumique d'acier $\phi$ (matériau composite, comme une \og nappe d'armatures \fg{} d'un
calcul aux éléments finis). L'aire d'acier par mètre de profondeur est $A_s/b$ ; pour un espacement
de particules $\Delta_p = \Delta x/2 = 7.8$~mm,
\[
  \phi = \frac{A_s/b}{\Delta_p}.
\]
Les lits du type 4 (tableau~2 de l'article) donnent ($b = 0.2$~m) :

\begin{center}
\begin{tabular}{@{}lrrrr@{}}
\toprule
lit & hauteur (mm) & $A_s$ (mm$^2$) & rangée de particules & $\phi$ \\
\midrule
4 (bas) & 19 & $4\times113.1$ & 2 & 0.290 \\
3 & 72 & $2\times56.5$ & 9 & 0.072 \\
2 & 128 & $2\times56.5$ & 16 & 0.072 \\
1 (haut) & 181 & $4\times113.1$ & 23 & 0.290 \\
\bottomrule
\end{tabular}
\end{center}

Sur la particule armée, la contrainte est la moyenne pondérée entre béton et barre :
\begin{equation}
  \bm\tau = (1-\phi)\,\bm\tau_{\rm béton} + \phi\,\sigma_s\ \mathbf n\otimes\mathbf n,
  \qquad \mathbf n = \frac{\mathbf F\mathbf e_x}{|\mathbf F\mathbf e_x|},\qquad
  \varepsilon_s = \ln|\mathbf F\mathbf e_x| .
\end{equation}
La barre suit la direction $x$ du matériau (elle tourne avec $\mathbf F$). Sa loi est élastique
parfaitement plastique ; la plasticité utilise un retour radial avec une variable d'état
$\varepsilon_p$ :
\begin{equation}
  \sigma_s^{\rm essai} = E_s(\varepsilon_s - \varepsilon_p),\quad
  \sigma_s = \operatorname{clamp}(\sigma_s^{\rm essai},\,-f_y,\,f_y),\quad
  \varepsilon_p \leftarrow \varepsilon_p + \frac{\sigma_s^{\rm essai} - \sigma_s}{E_s}.
\end{equation}
Les barres sont supposées parfaitement adhérentes : elles partagent la vitesse nodale du béton.

\subsection{Résistances aléatoires}

Chaque particule reçoit ses propres $f_{cr}$ et $f_c$, lognormales de moyenne celle du
tableau~3 de l'article et de coefficients de variation $0.26$ ($0.9/3.5$) et $0.11$ ($5.1/47.1$) :
\[
  X = m\,\exp\!\big(\sigma_{\ln}\,z - \tfrac12\sigma_{\ln}^2\big),\quad z\sim\mathcal N(0,1),\quad
  \sigma_{\ln} = \sqrt{\ln(1 + c_v^2)},
\]
tronquée pour éviter les valeurs aberrantes. Le tirage est reproductible (\code{random_seed=1}).
Il casse la symétrie parfaite : sans lui, toutes les particules d'une même fibre fissureraient au
même pas.

% ================================================================
\section{Pièges numériques}

\subsection{Unités et pas de temps}

Tout est en unités SI (m, kg, s, Pa). La poutre de $2.4$~m est placée dans un domaine de
$3\times0.625$~m ($192\times40$ n\oe uds, $\Delta x = 1.56$~cm), soit $7\,982$ particules
($307\times26$). Le pas de temps explicite doit tenir compte de la raideur \emph{axiale}
maximale, béton plus acier :
\[
  c_{\max} = \sqrt{\frac{\lambda + 2\mu + \phi_{\max} E_s}{\rho}} \approx 6.5\ \mathrm{km/s}
  \quad(c_p \text{ du béton seul : } 3.75\ \mathrm{km/s}),\qquad
  \Delta t = 0.3\,\frac{\Delta x}{c_{\max}} \approx 7\cdot10^{-7}\ \mathrm{s}.
\]

\subsection{Précision en \code{float32}}

Une déformation de $10^{-4}$ (fissuration) se joue à la quatrième décimale de $\mathbf F$, et un
pas de temps de $7\cdot10^{-7}$~s ajoute à $\mathbf F$ des incréments de l'ordre de $10^{-8}$ à $10^{-7}$,
soit la précision relative de \code{float32} autour de $1$ ($\approx 6\cdot10^{-8}$). Le code stocke donc
$\mathbf F_m = \mathbf F - \mathbf I$ (petit, donc précis) et le déplacement $\mathbf u = \mathbf x - \mathbf X_0$ plutôt
que les positions ; on reconstruit $\mathbf F$ et $\mathbf x$ au moment de s'en servir :
\begin{lstlisting}
Fm[p] += dt * C_new @ (I + Fm[p])       # dF/dt = (grad v) F  <=>  Fm <- Fm + dt C (I + Fm)
Up[p] += dt * v_s[p]                    # x = X0 + Up
\end{lstlisting}
C'est un raisonnement de précaution : nous n'avons pas mesuré l'erreur qu'on obtiendrait sans cela.

\subsection{Appuis et vérins glissants}

Une première version imposait une vitesse \emph{nulle} aux n\oe uds des appuis et une vitesse
imposée aux n\oe uds des vérins, dans toutes les directions. La poutre était alors partiellement
encastrée : la raideur mesurée valait environ $270$~kN/mm, contre environ $43$~kN/mm avec des
contacts glissants, où seule la composante normale est contrainte :
\begin{lstlisting}
if mask[i, j] == 1 and grid_v[i, j].y < 0:            # appui : la poutre ne peut pas descendre
    grid_v[i, j].y = 0.0
...
if in_plateau and grid_v[i, j].y > -v_load:           # verin : plateau qui descend a v_load
    react[None] += grid_m[i, j] * (grid_v[i, j].y + v_load) / dt
    grid_v[i, j].y = -v_load
\end{lstlisting}
La force que le vérin exerce est la quantité de mouvement qu'il doit retirer aux n\oe uds pour
imposer sa vitesse, $\sum_i m_i\,(v_{i,y} + v_{\rm load})/\Delta t$, accumulée dans \code{react}
et moyennée sur les pas d'une image. Elle est \emph{par mètre de profondeur} ; on multiplie par
la largeur $b = 0.2$~m pour obtenir des kN.

\subsection{Essai quasi statique}

On pilote les vérins en déplacement (vitesse $0.2$~m/s par défaut, soit $25$~mm en $\approx0.13$~s),
avec une montée en vitesse en \emph{smoothstep} sur $20$~ms, pour ne pas donner un choc à la poutre.
D'après une estimation d'ordre de grandeur, l'énergie cinétique de la poutre devient petite devant
l'énergie de déformation une fois la poutre chargée ; au début de l'essai, l'inertie fausse un peu la
raideur (voir les résultats). Un
amortissement de grille de $100$~s$^{-1}$ éteint les vibrations.

% ================================================================
\section{Le code, pas à pas}

Le programme complet (\code{MpM/beton.py}) est en annexe. On en commente ici les morceaux
qui portent la physique.

\subsection{Lois sécantes}

\begin{lstlisting}
@ti.func
def r_tension(kap, fcr_i):
    r = 1.0
    if kap > fcr_i / E_c:
        r = ti.min(1.0, fcr_i / (1.0 + ti.sqrt(500.0 * kap)) / (E_c * kap))
    return ti.max(r, r_min)

@ti.func
def r_compression(kap, fc_i):
    x = kap / (1.8 * fc_i / E_c)
    r = ti.min(1.0, (2.0 / 1.8) / (1.0 + x * x))
    if kap > eps_cu:
        t = ti.min((kap - eps_cu) / (0.5 * eps_cu), 1.0)
        r = r * (1.0 - t) + r_crush * t
    return ti.max(r, r_min)
\end{lstlisting}

\subsection{Contrainte de Kirchhoff}

\begin{lstlisting}
@ti.func
def stress_kirchhoff(F, kt, kc, ph, ep, fcr_i, fc_i, nl: int):
    U, sig, V = ti.svd(F)
    e1 = ti.log(ti.max(sig[0, 0], 0.05))
    e2 = ti.log(ti.max(sig[1, 1], 0.05))
    tr = e1 + e2
    s1 = 2.0 * mu_c * e1 + la_c * tr                 # contraintes elastiques principales
    s2 = 2.0 * mu_c * e2 + la_c * tr
    if nl == 1:                                      # facteur selon le signe (unilateral)
        rt = r_tension(kt, fcr_i)
        rc = r_compression(kc, fc_i)
        if s1 > 0.0: s1 *= rt
        else:        s1 *= rc
        if s2 > 0.0: s2 *= rt
        else:        s2 *= rc
    S = ti.Matrix([[s1, 0.0], [0.0, s2]])
    tau = U @ S @ U.transpose()
    if ph > 0.0:                                     # particule armee : barre uniaxiale
        n = ti.Vector([F[0, 0], F[1, 0]])
        lam = n.norm()
        n = n / lam
        ss = E_s * ti.log(lam)
        if nl == 1:
            ss = ti.math.clamp(E_s * (ti.log(lam) - ep), -fy, fy)
        tau = (1.0 - ph) * tau + ph * ss * n.outer_product(n)
    return tau
\end{lstlisting}
Le paramètre \code{nl} bascule entre le modèle \emph{élastique linéaire} (touche 1) et le modèle
\emph{béton armé non linéaire} (touche 2).

\subsection{Lissage et historique}

\code{scatter_eps} calcule $\varepsilon_t$ et $\varepsilon_c$ par particule et les dépose sur la grille ;
\code{update_state} les relit, met à jour $\kappa_t$, $\kappa_c$ (jamais décroissants) et la plasticité de
l'acier :
\begin{lstlisting}
kap_t[p] = ti.max(kap_t[p], et)
kap_c[p] = ti.max(kap_c[p], ec)
if phi[p] > 0.0:
    eps = ti.log(ti.Vector([Fm[p][0, 0] + 1.0, Fm[p][1, 0]]).norm())
    trial = E_s * (eps - eps_p[p])
    ss = ti.math.clamp(trial, -fy, fy)
    eps_p[p] += (trial - ss) / E_s                    # retour radial
\end{lstlisting}

\subsection{Un pas de temps}

\begin{lstlisting}
def substep(v_load, model):
    s = smoothstep(State.t)                 # montee en charge progressive
    v_eff = v_load * s
    nl = 1 if model >= 2 else 0
    clear_grid()
    P2G_solid(dt, nl)
    grid_update(dt, g * s, v_eff, State.u)  # appuis, verins, reaction
    G2P_solid(dt)
    if nl == 1:
        scatter_eps()
        update_state()
    State.t += dt
    State.u += v_eff * dt                   # course des verins
\end{lstlisting}

% ================================================================
\section{Résultats}

\subsection{Courbe charge--flèche}

\begin{center}
\includegraphics[width=0.78\linewidth]{tuto_beton_courbe.png}
\end{center}

La courbe en tirets est l'analyse de section de \code{MpM/beton_section.py}, qui applique les mêmes lois
à la manière de l'article (couches de béton, quatre lits d'acier, courbure intégrée le long de la
poutre) et s'arrête au pic du moment. Chiffres relevés (charge $P$ totale en kN) :

\begin{center}
\begin{tabular}{@{}lrrrrrrr@{}}
\toprule
flèche à mi-portée (mm) & 0.5 & 1 & 2 & 3 & 5 & 10 & 15 \\
\midrule
MPM & 29.1 & 44.1 & 65.1 & 84.7 & 124.5 & 131.0 & 123.7 \\
analyse de section & 20.8 & 33.3 & 51.5 & 68.1 & 101.0 & 110.5 & 113.1 \\
\bottomrule
\end{tabular}
\end{center}

Trois régimes apparaissent, comme sur les courbes de l'article : une phase élastique
(pente initiale de l'analyse de section : $41.5$~kN/mm), une phase fissurée de raideur réduite, où
l'acier prend le relais du béton tendu, puis un \emph{plateau} de plastification des aciers
tendus. Dans l'analyse de section, le plateau est à $113$~kN (moment $34$~kN\,m) ; le MPM le
dépasse d'environ $16$~\% ($\approx131$~kN). À $25$~mm de course de vérin, aucune particule n'a
atteint $\varepsilon_{cu}$ : le béton comprimé n'est pas encore écrasé, donc on ne voit pas la
ruine, seulement le début de l'adoucissement (charge autour de $115$ à $120$~kN).

\paragraph{Écarts et ce qu'on peut en dire.} Le MPM est plus raide au début (pente $\approx 54$~kN/mm
mesurée entre $0.15$ et $0.45$~mm, pour $41.5$ en section). Nous avons vérifié qu'à vitesse de
vérin plus faible ($0.05$~m/s), avec le modèle élastique linéaire, la pente descend à
$38$--$48$~kN/mm : une partie de l'écart vient donc de l'inertie et de la montée en charge. Les
autres écarts (plateau plus haut de $16$~\%, transition plus tardive) ne sont pas élucidés ici ;
pistes plausibles, non testées : la loi de raidissement~\eqref{eq:tension} appliquée à tout le
béton tendu (dans l'article, elle ne concerne que la zone proche des barres), les déformations planes
($E/(1-\nu^2)$), et la modélisation d'une barre par une simple rangée de particules. Notons que
l'article observe aussi une surestimation systématique de la capacité portante. Un pic isolé de
$138$~kN vers $17$~mm est un artefact numérique (une seule série de pas).

\subsection{Cartes d'endommagement et de contraintes}

Couleurs du mode \og endommagement \fg{} : gris (intact) $\to$ jaune (béton fissuré) $\to$ rouge
($d_t = 1 - r_t > 0.95$), cyan si la compression est endommagée, violet si écrasé ; les lits
d'acier sont bleus, et oranges lorsqu'ils ont plastifié.

\begin{center}
\includegraphics[width=\linewidth]{tuto_beton_endo_a.png}\\[2pt]
\includegraphics[width=\linewidth]{tuto_beton_endo_b.png}
\end{center}

En haut, course de vérin de $4$~mm (flèche $5.1$~mm, $P = 125$~kN) : les premières fissures de
flexion (rouge) sont dans la fibre inférieure, sur toute la zone de moment constant, avec une
bande jaune au-dessus. En bas, course de $14$~mm (flèche $16.4$~mm, $P = 118$~kN) : la zone
fissurée a gagné presque toute la hauteur au centre, avec des limites obliques dans les travées
d'effort tranchant, et les lits d'armatures inférieurs sont orange (plastifiés) dans la zone centrale. La
fissuration est \emph{diffuse} : on voit une zone endommagée épaisse, pas des fissures isolées
(voir les limites).

Le mode \og contrainte principale \fg{} (touche Espace) montre la traction en rouge sous
l'axe neutre et la compression en bleu au-dessus, le diagramme de flexion classique :

\begin{center}
\includegraphics[width=\linewidth]{tuto_beton_contrainte_a.png}
\end{center}

% ================================================================
\section{Ce qui n'a pas marché}

\paragraph{Rupture cohésive : instable.} Le tutoriel précédent faisait passer une particule
totalement endommagée à l'état \og rompu \fg{} (perte de cohésion, allongements principaux écrêtés à $1$).
Appliqué ici (rupture dès que $\kappa_t \ge 1.5\cdot10^{-3}$ pour le béton non armé, écrasement
à $\kappa_c \ge \varepsilon_{cu}$), le calcul se déroulait normalement jusqu'à $P \approx 112$~kN, puis
la fissuration se propageait sur toute la hauteur en quelques pas : la vitesse maximale des
particules passait de $0.14$ à plus de $50$~m/s, la flèche de $4.5$ à $50$~mm en $5$~ms et la
poutre était pulvérisée. Ce n'est pas une rupture de barre : en désactivant la rupture des
barres, l'effondrement était identique. En revanche, le modèle sans rupture cohésive (endommagement
unilatéral seul, celui décrit plus haut) est stable sur toute la course. Le mécanisme de
la rupture cohésive, pertinent pour un solide qui doit se \emph{fragmenter}, est donc retiré ici.
La cause exacte de l'instabilité n'a pas été établie ; l'hypothèse la plus probable est la perte
brutale des $\approx 20$~kN que fournissait le béton tendu (raidissement) au moment où la fissure traverse
la zone tendue, l'énergie libérée étant amplifiée par l'écrêtage de $\mathbf F$.

\paragraph{Appuis collants.} Voir plus haut : $270$~kN/mm au lieu de $43$~kN/mm.

\paragraph{Rupture des barres à $2$~\%.} Nous n'avons pas implémenté la rupture des barres à
$\varepsilon = 0.02$. Cette valeur de l'article est une déformation \emph{moyenne} d'une section, alors
que dans une simulation où la fissure se localise sur une ou deux colonnes de particules, la
déformation de la barre y dépasse $2$~\% presque immédiatement. Il faudrait la lisser sur une longueur
d'adhérence ; ce n'est pas fait, et sans elle, la poutre ne \og casse \fg{} qu'en compression.

\paragraph{Rate limiting.} Le paramètre $\tau_D$ (vitesse d'endommagement limitée) du tutoriel précédent
n'est pas repris : l'essai étant piloté en déplacement, il n'a pas été nécessaire. Ce n'est pas
démontré pour un chargement dynamique.

% ================================================================
\section{Limites et pour aller plus loin}

\begin{enumerate}
  \item \textbf{Fissuration diffuse.} Le modèle est un modèle de fissures \emph{réparties} (comme
        l'article, qui calcule des courbures moyennes par section). Pour des fissures discrètes,
        il faudrait un critère de rupture avec énergie de fissuration régularisée, ou une
        méthode dédiée (CPIC, réseau de liaisons).
  \item \textbf{Ruine.} Aucune particule n'a atteint $\varepsilon_{cu}$ sur les $25$~mm simulés : la ruine
        par écrasement du béton comprimé n'a pas été observée. Il faut prolonger la course des vérins
        (le comportement après $\varepsilon_{cu}$ n'est donc pas validé).
  \item \textbf{Poutres précontraintes.} L'article traite aussi la précontrainte (pré- et
        post-tension, loi empirique de Collins et Mitchell). Elle s'ajouterait comme une déformation
        initiale dans les rangées d'acier.
  \item \textbf{Cycles et écrouissage.} Non pris en compte ici ni dans l'article.
  \item \textbf{2D en déformations planes.} Une barre est une rangée de particules, non un cylindre ;
        le confinement des étriers et l'adhérence acier--béton (glissement) sont ignorés.
  \item \textbf{Monte-Carlo.} On tire une réalisation. Pour retrouver une moyenne et un écart type de
        capacité portante comme dans l'article, il faut lancer beaucoup de tirages (\code{random_seed} variable) ;
        chaque essai coûte quelques minutes.
\end{enumerate}

\paragraph{Expériences.} Pour vérifier que chaque ingrédient joue son rôle, essayez :
\begin{itemize}
  \item mettre \code{cv_fcr = 0} : les fissures de la fibre inférieure naissent plus régulièrement ;
  \item supprimer les armatures (\code{phi_np[:] = 0}) : la poutre n'a plus de plateau et perd sa
        capacité peu après la fissuration ;
  \item passer à $f_{cr} = 7$~MPa ou à $E_c$ différent ;
  \item augmenter la vitesse des vérins (curseur) pour voir apparaître les effets dynamiques ;
  \item comparer touche~1 (élastique) et touche~2 (béton armé) à la même charge.
\end{itemize}

% ================================================================
\appendix
\section{Code complet : \texttt{MpM/beton.py}}

\begin{lstlisting}
# MpM/beton.py -- Poutre en béton armé en flexion 4 points, MLS-MPM 2D
#
# Inspiré de Cremona & Houde, "Modélisation déterministe et probabiliste du comportement
# mécanique simplifié des corps d'épreuve" (poutres de la Rance) :
#   - béton en compression : loi de Desayi/Neville, écrasement vers -0.004
#   - béton en traction    : linéaire jusqu'à fcr, puis fcr / (1 + sqrt(500 eps))  (raidissement)
#   - acier                : bilinéaire (élastique parfaitement plastique)
#   - résistances aléatoires par particule (lois lognormales, comme dans l'approche Monte-Carlo)
#   - essai de flexion 4 points sur appuis simples, portée 2 m, charges à 0.6 m des appuis
#
# Unités SI (m, kg, s, Pa). Le calcul est 2D en déformations planes : les efforts sont par mètre de
# profondeur ; on les multiplie par la largeur de la poutre b = 0.2 m pour retrouver des kN.
#
# Lancer :  python MpM/beton.py                 (fenêtre interactive)
#           python MpM/beton.py --headless 25   (course de vérin 25 mm, écrit beton_courbe.csv)

import sys
import numpy as np
import taichi as ti

ti.init(arch=ti.gpu, random_seed=1)

# ---------------------------------------------------------------- domaine (rectangle Lx x Ly)
Lx, Ly = 3.0, 0.625
nx, ny = 192, 40
dx = Lx / nx                          # 1.56 cm
inv_dx = 1.0 / dx
bound = 3
render_substep = 50

# ---------------------------------------------------------------- béton (poutre 421 de la Rance)
rho_c = 2400.0
E_c = 30.4e9                          # module d'Young
nu_c = 0.2
fc = 47.1e6                           # résistance en compression
fcr = 3.5e6                           # résistance en traction
cv_fc, cv_fcr = 0.11, 0.26            # coefficients de variation (écart type / moyenne, tableau 3)
eps_cu = 0.004                        # déformation d'écrasement
r_min = 1e-3                          # raideur résiduelle en traction (béton fissuré)
r_crush = 0.1                         # raideur résiduelle en compression (béton écrasé)
mu_c = E_c / (2 * (1 + nu_c))
la_c = E_c * nu_c / ((1 + nu_c) * (1 - 2 * nu_c))

# ---------------------------------------------------------------- acier d'armature
E_s = 230e9
fy = 309e6                            # limite d'élasticité

# ---------------------------------------------------------------- amortissement et chargement
damp = 100.0                          # amortissement de la vitesse de grille (1/s)
g = 9.81
t_ramp = 0.02                         # durée de la montée en charge (s)
b_depth = 0.2                         # largeur de la poutre (m) : passage N/m -> N

# ---------------------------------------------------------------- géométrie
beam_x0, beam_x1 = 0.3, 2.7           # la poutre fait 2.4 m
beam_y0 = 0.25
xs1, xs2 = 0.5, 2.5                   # appuis : portée 2.0 m
xp1, xp2 = 1.1, 1.9                   # charges à 0.6 m des appuis
hw = 0.05                             # demi-largeur des appuis et des vérins
h_target = 0.2

p_spacing = 0.5 * dx                  # 2 x 2 particules par cellule
p_vol = p_spacing ** 2
p_mass = p_vol * rho_c
n_px = int(round((beam_x1 - beam_x0) / p_spacing))
n_py = int(round(h_target / p_spacing))
n_solid = n_px * n_py
h_beam = n_py * p_spacing
beam_y1 = beam_y0 + h_beam
x_mid = 0.5 * (xs1 + xs2)

# lits d'armatures du type 4 de l'article : (distance à la base en m, aire totale en m2)
rebars = [(0.019, 4 * 113.1e-6), (0.072, 2 * 56.5e-6), (0.128, 2 * 56.5e-6), (0.181, 4 * 113.1e-6)]
phi_np = np.zeros(n_py, dtype=np.float32)
for dist, area in rebars:
    phi_np[int(dist / p_spacing)] += (area / b_depth) / p_spacing      # fraction volumique d'acier
phi_max = float(phi_np.max())

# pas de temps explicite : dt < dx / c avec la raideur axiale maximale (béton + acier)
c_max = ((la_c + 2 * mu_c + phi_max * E_s) / rho_c) ** 0.5
dt = 0.3 * dx / c_max

sl_cr = float(np.log(1 + cv_fcr ** 2) ** 0.5)      # écarts types des logarithmes (lognormale)
sl_fc = float(np.log(1 + cv_fc ** 2) ** 0.5)

# ---------------------------------------------------------------- champs
X0 = ti.Vector.field(2, ti.f32, n_solid)          # position initiale
Up = ti.Vector.field(2, ti.f32, n_solid)          # déplacement (x = X0 + Up, voir précision float32)
v_s = ti.Vector.field(2, ti.f32, n_solid)
C_s = ti.Matrix.field(2, 2, ti.f32, n_solid)      # matrice affine APIC
Fm = ti.Matrix.field(2, 2, ti.f32, n_solid)       # F - I (l'écart à l'identité garde la précision)
kap_t = ti.field(ti.f32, n_solid)                 # plus grande déformation de traction équivalente
kap_c = ti.field(ti.f32, n_solid)                 # plus grande déformation de compression équivalente
phi = ti.field(ti.f32, n_solid)                   # fraction volumique d'acier de la particule
eps_p = ti.field(ti.f32, n_solid)                 # déformation plastique de l'acier
fcr_p = ti.field(ti.f32, n_solid)                 # résistance en traction de la particule
fc_p = ti.field(ti.f32, n_solid)                  # résistance en compression de la particule
phi_row = ti.field(ti.f32, n_py)
phi_row.from_numpy(phi_np)

grid_v = ti.Vector.field(2, ti.f32, (nx, ny))
grid_m = ti.field(ti.f32, (nx, ny))
grid_et = ti.field(ti.f32, (nx, ny))              # déformation de traction lissée (non local)
grid_ec = ti.field(ti.f32, (nx, ny))
grid_w = ti.field(ti.f32, (nx, ny))
mask = ti.field(ti.i32, (nx, ny))                 # 1 = appui
react = ti.field(ti.f32, ())                      # réaction des vérins (N/m)

W, H = 1200, 250
Image = ti.Vector.field(3, ti.f32, (W, H))


# ---------------------------------------------------------------- lois de comportement
@ti.func
def r_tension(kap, fcr_i):
    """Raideur sécante en traction : 1 avant fissuration, puis fcr/(1+sqrt(500 eps)) / (E eps)."""
    r = 1.0
    if kap > fcr_i / E_c:
        r = ti.min(1.0, fcr_i / (1.0 + ti.sqrt(500.0 * kap)) / (E_c * kap))
    return ti.max(r, r_min)


@ti.func
def r_compression(kap, fc_i):
    """Raideur sécante en compression : loi de Desayi, puis écrasement progressif après eps_cu."""
    x = kap / (1.8 * fc_i / E_c)
    r = ti.min(1.0, (2.0 / 1.8) / (1.0 + x * x))
    if kap > eps_cu:
        t = ti.min((kap - eps_cu) / (0.5 * eps_cu), 1.0)
        r = r * (1.0 - t) + r_crush * t
    return ti.max(r, r_min)


@ti.func
def stress_kirchhoff(F, kt, kc, ph, ep, fcr_i, fc_i, nl: int):
    """Contrainte de Kirchhoff : béton (Hencky, raideurs sécantes par direction principale)
    + acier (barre uniaxiale suivant la direction x du matériau, élastique parfaitement plastique)."""
    U, sig, V = ti.svd(F)
    e1 = ti.log(ti.max(sig[0, 0], 0.05))
    e2 = ti.log(ti.max(sig[1, 1], 0.05))
    tr = e1 + e2
    s1 = 2.0 * mu_c * e1 + la_c * tr
    s2 = 2.0 * mu_c * e2 + la_c * tr
    if nl == 1:
        rt = r_tension(kt, fcr_i)
        rc = r_compression(kc, fc_i)
        if s1 > 0.0:
            s1 *= rt
        else:
            s1 *= rc
        if s2 > 0.0:
            s2 *= rt
        else:
            s2 *= rc
    S = ti.Matrix([[s1, 0.0], [0.0, s2]])
    tau = U @ S @ U.transpose()
    if ph > 0.0:
        n = ti.Vector([F[0, 0], F[1, 0]])
        lam = n.norm()
        n = n / lam
        ss = E_s * ti.log(lam)
        if nl == 1:
            ss = ti.math.clamp(E_s * (ti.log(lam) - ep), -fy, fy)
        tau = (1.0 - ph) * tau + ph * ss * n.outer_product(n)
    return tau


# ---------------------------------------------------------------- 1. particules -> grille
@ti.kernel
def P2G_solid(dt: float, nl: int):
    for p in X0:
        xp = X0[p] + Up[p]
        base = (xp * inv_dx - 0.5).cast(int)
        fx = xp * inv_dx - base
        w = [0.5 * (1.5 - fx)**2, 0.75 - (fx - 1.0)**2, 0.5 * (fx - 0.5)**2]

        F = ti.Matrix.identity(ti.f32, 2) + Fm[p]
        tau = stress_kirchhoff(F, kap_t[p], kap_c[p], phi[p], eps_p[p], fcr_p[p], fc_p[p], nl)
        affine = (-dt * p_vol * 4.0 * inv_dx * inv_dx) * tau + p_mass * C_s[p]

        for i, j in ti.static(ti.ndrange(3, 3)):
            offset = ti.Vector([i, j])
            d_pos = (offset - fx) * dx
            weight = w[i].x * w[j].y
            grid_m[base + offset] += weight * p_mass
            grid_v[base + offset] += weight * (p_mass * v_s[p] + affine @ d_pos)


# ---------------------------------------------------------------- 2. grille : appuis, vérins
@ti.kernel
def grid_update(dt: float, g_eff: float, v_load: float, u_load: float):
    for i, j in grid_m:
        if grid_m[i, j] > 0:
            grid_v[i, j] /= grid_m[i, j]
            grid_v[i, j] *= 1.0 - damp * dt
            grid_v[i, j].y -= dt * g_eff

            if i < bound and grid_v[i, j].x < 0:
                grid_v[i, j].x = 0.0
            if i > nx - bound and grid_v[i, j].x > 0:
                grid_v[i, j].x = 0.0
            if j < bound and grid_v[i, j].y < 0:
                grid_v[i, j].y = 0.0
            if j > ny - bound and grid_v[i, j].y > 0:
                grid_v[i, j].y = 0.0

            # appuis : contact glissant sans frottement (seule la composante normale est contrainte)
            if mask[i, j] == 1 and grid_v[i, j].y < 0:
                grid_v[i, j].y = 0.0

            # vérins : plateaux glissants qui descendent à la vitesse v_load (composante normale imposée),
            # noeuds situés sous la face inférieure des plateaux
            pos = ti.Vector([i, j]) * dx
            if (ti.abs(pos.x - xp1) <= hw or ti.abs(pos.x - xp2) <= hw) and pos.y >= beam_y1 - u_load:
                if grid_v[i, j].y > -v_load:
                    react[None] += grid_m[i, j] * (grid_v[i, j].y + v_load) / dt    # force sur la poutre
                    grid_v[i, j].y = -v_load


@ti.kernel
def init_mask():
    for i, j in mask:
        pos = ti.Vector([i, j]) * dx
        support = (ti.abs(pos.x - xs1) <= hw or ti.abs(pos.x - xs2) <= hw) and pos.y <= beam_y0
        mask[i, j] = 1 if support else 0


# ---------------------------------------------------------------- 3. grille -> particules
@ti.kernel
def G2P_solid(dt: float):
    I = ti.Matrix.identity(ti.f32, 2)
    for p in X0:
        xp = X0[p] + Up[p]
        base = (xp * inv_dx - 0.5).cast(int)
        fx = xp * inv_dx - base
        w = [0.5 * (1.5 - fx)**2, 0.75 - (fx - 1.0)**2, 0.5 * (fx - 0.5)**2]

        v_new = ti.Vector.zero(ti.f32, 2)
        C_new = ti.Matrix.zero(ti.f32, 2, 2)
        for i, j in ti.static(ti.ndrange(3, 3)):
            offset = ti.Vector([i, j])
            d_pos = (offset - fx) * dx
            weight = w[i].x * w[j].y
            v_new += weight * grid_v[base + offset]
            C_new += weight * grid_v[base + offset].outer_product(d_pos) * (4.0 * inv_dx * inv_dx)
        v_s[p] = v_new
        C_s[p] = C_new

        # F = I + Fm, dF/dt = (grad v) F   ->   Fm <- Fm + dt C (I + Fm)
        Fm[p] += dt * C_new @ (I + Fm[p])

        Up[p] += dt * v_s[p]
        xn = X0[p] + Up[p]
        xc = ti.math.clamp(xn, ti.Vector([bound * dx, bound * dx]),
                           ti.Vector([Lx - bound * dx, Ly - bound * dx]))
        Up[p] += xc - xn


# ---------------------------------------------------------------- 4. endommagement non local
@ti.kernel
def clear_grid():
    for i, j in grid_m:
        grid_v[i, j] = [0.0, 0.0]
        grid_m[i, j] = 0.0
        grid_et[i, j] = 0.0
        grid_ec[i, j] = 0.0
        grid_w[i, j] = 0.0


@ti.kernel
def scatter_eps():
    """Étale sur la grille les déformations équivalentes (contrainte principale élastique / E)."""
    for p in X0:
        xp = X0[p] + Up[p]
        base = (xp * inv_dx - 0.5).cast(int)
        fx = xp * inv_dx - base
        w = [0.5 * (1.5 - fx)**2, 0.75 - (fx - 1.0)**2, 0.5 * (fx - 0.5)**2]

        Us, sg, Vs = ti.svd(ti.Matrix.identity(ti.f32, 2) + Fm[p])
        e1 = ti.log(ti.max(sg[0, 0], 0.05))
        e2 = ti.log(ti.max(sg[1, 1], 0.05))
        s1 = 2.0 * mu_c * e1 + la_c * (e1 + e2)
        s2 = 2.0 * mu_c * e2 + la_c * (e1 + e2)
        et = ti.max(ti.max(s1, s2), 0.0) / E_c
        ec = ti.max(ti.max(-s1, -s2), 0.0) / E_c

        for i, j in ti.static(ti.ndrange(3, 3)):
            weight = w[i].x * w[j].y
            node = base + ti.Vector([i, j])
            grid_et[node] += weight * et
            grid_ec[node] += weight * ec
            grid_w[node] += weight


@ti.kernel
def update_state():
    """Historique de déformation (irréversible) du béton, plasticité de l'acier."""
    for p in X0:
        xp = X0[p] + Up[p]
        base = (xp * inv_dx - 0.5).cast(int)
        fx = xp * inv_dx - base
        w = [0.5 * (1.5 - fx)**2, 0.75 - (fx - 1.0)**2, 0.5 * (fx - 0.5)**2]

        et = 0.0
        ec = 0.0
        for i, j in ti.static(ti.ndrange(3, 3)):
            node = base + ti.Vector([i, j])
            if grid_w[node] > 0:
                weight = w[i].x * w[j].y
                et += weight * grid_et[node] / grid_w[node]
                ec += weight * grid_ec[node] / grid_w[node]
        kap_t[p] = ti.max(kap_t[p], et)
        kap_c[p] = ti.max(kap_c[p], ec)

        if phi[p] > 0.0:
            eps = ti.log(ti.Vector([Fm[p][0, 0] + 1.0, Fm[p][1, 0]]).norm())
            trial = E_s * (eps - eps_p[p])
            ss = ti.math.clamp(trial, -fy, fy)
            eps_p[p] += (trial - ss) / E_s                  # retour radial : écoulement plastique


# ---------------------------------------------------------------- initialisation, mesures
@ti.kernel
def init_beam():
    for p in X0:
        i = p % n_px
        j = p // n_px
        X0[p] = [beam_x0 + (i + 0.5) * p_spacing, beam_y0 + (j + 0.5) * p_spacing]
        Up[p] = [0.0, 0.0]
        v_s[p] = [0.0, 0.0]
        C_s[p] = ti.Matrix.zero(ti.f32, 2, 2)
        Fm[p] = ti.Matrix.zero(ti.f32, 2, 2)
        kap_t[p] = 0.0
        kap_c[p] = 0.0
        eps_p[p] = 0.0
        phi[p] = phi_row[j]
        # résistances aléatoires par particule : lognormales de moyenne fcr, fc
        fcr_p[p] = fcr * ti.min(ti.max(ti.exp(sl_cr * ti.randn() - 0.5 * sl_cr**2), 0.3), 3.0)
        fc_p[p] = fc * ti.min(ti.max(ti.exp(sl_fc * ti.randn() - 0.5 * sl_fc**2), 0.5), 2.0)


@ti.kernel
def bottom_mid_uy() -> ti.types.vector(2, ti.f32):
    """Somme et nombre des déplacements verticaux de la fibre inférieure à mi-portée."""
    s = 0.0
    n = 0.0
    for p in X0:
        if p // n_px == 0 and ti.abs(X0[p].x - x_mid) < 0.1:
            s += Up[p].y
            n += 1.0
    return ti.Vector([s, n])


@ti.kernel
def damage_counts() -> ti.types.vector(2, ti.f32):
    """Nombre de particules de béton fissurées (d_t > 0.9) et écrasées (eps_c > eps_cu)."""
    nt = 0.0
    nc = 0.0
    for p in X0:
        if phi[p] == 0.0:
            if 1.0 - r_tension(kap_t[p], fcr_p[p]) > 0.9:
                nt += 1.0
            if kap_c[p] > eps_cu:
                nc += 1.0
    return ti.Vector([nt, nc])


# ---------------------------------------------------------------- affichage
@ti.func
def particle_color(p: int, mode: int):
    col = ti.Vector([0.75, 0.75, 0.75])
    if mode == 0:
        if phi[p] > 0.0:
            col = ti.Vector([0.25, 0.45, 1.0])                           # armature élastique
            if ti.abs(eps_p[p]) > 1e-4:
                col = ti.Vector([1.0, 0.55, 0.10])                       # armature plastifiée
        else:
            d_t = 1.0 - r_tension(kap_t[p], fcr_p[p])
            d_c = ti.min(1.0, (1.0 - r_compression(kap_c[p], fc_p[p])) / 0.6)
            col = ti.Vector([0.75, 0.75, 0.75]) * (1 - d_t) + ti.Vector([1.0, 0.85, 0.1]) * d_t
            if d_t > 0.95:
                col = ti.Vector([0.90, 0.15, 0.15])                      # fissure
            col = col * (1 - d_c) + ti.Vector([0.3, 0.85, 0.95]) * d_c   # compression endommagée
            if kap_c[p] > eps_cu:
                col = ti.Vector([0.60, 0.20, 0.80])                      # écrasement
    else:
        Us, sg, Vs = ti.svd(ti.Matrix.identity(ti.f32, 2) + Fm[p])
        e1 = ti.log(ti.max(sg[0, 0], 0.05))
        e2 = ti.log(ti.max(sg[1, 1], 0.05))
        s1 = 2.0 * mu_c * e1 + la_c * (e1 + e2)
        s2 = 2.0 * mu_c * e2 + la_c * (e1 + e2)
        a = ti.math.clamp(ti.max(s1, s2) / fcr, 0.0, 1.0)
        b = ti.math.clamp(-ti.min(s1, s2) / fc, 0.0, 1.0)
        t = 0.5 + 0.5 * (a - b)
        col = ti.Vector([0.2, 0.4, 1.0]) * (1 - t) + ti.Vector([1.0, 0.25, 0.1]) * t
    return col


@ti.kernel
def render(mode: int, u_load: float):
    for i, j in Image:
        pos = ti.Vector([(i + 0.5) * Lx / W, (j + 0.5) * Ly / H])
        col = ti.Vector([0.02, 0.02, 0.08])
        if (ti.abs(pos.x - xs1) <= hw or ti.abs(pos.x - xs2) <= hw) and pos.y <= beam_y0:
            col = ti.Vector([0.35, 0.35, 0.35])
        if (ti.abs(pos.x - xp1) <= hw or ti.abs(pos.x - xp2) <= hw) and pos.y >= beam_y1 - u_load:
            col = ti.Vector([0.55, 0.55, 0.60])
        Image[i, j] = col
    for p in X0:
        xp = X0[p] + Up[p]
        ci = int(xp.x / Lx * W)
        cj = int(xp.y / Ly * H)
        col = particle_color(p, mode)
        for a, b in ti.static(ti.ndrange(4, 4)):            # carré de 4 x 4 pixels (pas de trous)
            ii = ci + a - 1
            jj = cj + b - 1
            if 0 <= ii < W and 0 <= jj < H:
                Image[ii, jj] = col


# ---------------------------------------------------------------- simulation
class State:
    t = 0.0
    u = 0.0          # course des vérins (m)


def reset():
    init_beam()
    init_mask()
    State.t = 0.0
    State.u = 0.0


def smoothstep(t):
    s = min(t / t_ramp, 1.0)
    return s * s * (3.0 - 2.0 * s)


def substep(v_load, model):
    """Un pas de temps. model : 1 élastique linéaire, 2 béton armé non linéaire."""
    s = smoothstep(State.t)
    v_eff = v_load * s
    nl = 1 if model >= 2 else 0
    clear_grid()
    P2G_solid(dt, nl)
    grid_update(dt, g * s, v_eff, State.u)
    G2P_solid(dt)
    if nl == 1:
        scatter_eps()
        update_state()
    State.t += dt
    State.u += v_eff * dt


def advance(n, v_load, model):
    """n pas de temps ; retourne (P en kN moyenné sur la série, flèche à mi-portée en mm)."""
    react[None] = 0.0
    for _ in range(n):
        substep(v_load, model)
    P = react[None] / n * b_depth / 1000.0
    sn = bottom_mid_uy()
    delta = -sn[0] / max(sn[1], 1.0) * 1000.0
    return P, delta


HEADER = "t_s,course_verin_mm,fleche_mm,P_kN,n_fissurees,n_ecrasees"


def run_headless(u_max_mm, v_load=0.2, model=2, csv="beton_courbe.csv"):
    reset()
    rows = []
    while State.u * 1000.0 < u_max_mm:
        P, delta = advance(render_substep, v_load, model)
        dc = damage_counts()
        rows.append((State.t, State.u * 1000.0, delta, P, dc[0], dc[1]))
    np.savetxt(csv, np.array(rows), delimiter=",", header=HEADER, comments="")
    return np.array(rows)


def main():
    window = ti.ui.Window("MPM 2D - poutre en béton armé, flexion 4 points", res=(W, H))
    canvas = window.get_canvas()
    gui = window.get_gui()
    model = 2
    color_mode = 0
    v_load = 0.2
    P_max = 0.0
    reset()
    hist = []

    while window.running:
        while window.get_event(ti.ui.PRESS):
            k = window.event.key
            if k == ti.ui.ESCAPE:
                window.running = False
            elif k == ti.ui.SPACE:
                color_mode = 1 - color_mode
            elif k == 'r':
                reset()
                P_max = 0.0
                hist = []
            elif k in ('1', '2'):
                model = int(k)
            elif k == 's' and hist:
                np.savetxt("beton_courbe.csv", np.array(hist), delimiter=",", header=HEADER, comments="")

        P, delta = advance(render_substep, v_load, model)
        P_max = max(P_max, P)
        dc = damage_counts()
        hist.append((State.t, State.u * 1000.0, delta, P, dc[0], dc[1]))

        render(color_mode, State.u)
        canvas.set_image(Image)

        gui.begin("Flexion 4 points", 0.01, 0.01, 0.42, 0.40)
        v_load = gui.slider_float("vitesse des vérins (m/s)", v_load, 0.05, 0.5)
        gui.text(f"modèle : {['', 'élastique linéaire', 'béton armé non linéaire'][model]}   (touches 1/2)")
        gui.text(f"t = {State.t * 1000:.1f} ms   course = {State.u * 1000:.1f} mm")
        gui.text(f"charge P = {P:.1f} kN   (max {P_max:.1f} kN)")
        gui.text(f"flèche à mi-portée = {delta:.2f} mm")
        gui.text(f"béton fissuré = {int(dc[0])}   écrasé = {int(dc[1])}   (sur {n_solid})")
        gui.text("couleur : " + ("endommagement, acier" if color_mode == 0 else "contrainte principale") + "  (ESPACE)")
        gui.text("R : reset   S : sauver la courbe   ESC : quitter")
        gui.end()
        window.show()


if __name__ == "__main__":
    if len(sys.argv) >= 2 and sys.argv[1] == "--headless":
        umax = float(sys.argv[2]) if len(sys.argv) >= 3 else 25.0
        res = run_headless(umax)
        step = max(1, len(res) // 25)
        print(" t(ms)  course(mm)  fleche(mm)  P(kN)  fissurees  ecrasees")
        for r in res[::step]:
            print(f"{r[0]*1000:7.1f} {r[1]:9.2f} {r[2]:10.2f} {r[3]:8.1f} {int(r[4]):9d} {int(r[5]):9d}")
    else:
        main()
\end{lstlisting}

\section{Référence RDM : \texttt{MpM/beton\_section.py}}

\begin{lstlisting}
# MpM/beton_section.py -- Référence "résistance des matériaux" à comparer au calcul MPM
#
# Analyse de section à la manière de l'article (2.1) : pour une courbure phi donnée, on cherche la
# déformation au centre de gravité qui annule l'effort normal N, puis on calcule le moment M(phi).
# On intègre ensuite les courbures le long de la poutre (flexion 4 points) pour obtenir la flèche.
#
# Lancer :  python MpM/beton_section.py [beton_courbe.csv]   (trace la comparaison si le csv existe)

import sys
import numpy as np

# mêmes données que MpM/beton.py
b, h = 0.2, 0.2
E_c, fc, fcr, eps_cu = 30.4e9, 47.1e6, 3.5e6, 0.004
E_s, fy = 230e9, 309e6
rebars = [(0.019, 4 * 113.1e-6), (0.072, 2 * 56.5e-6), (0.128, 2 * 56.5e-6), (0.181, 4 * 113.1e-6)]
L, a = 2.0, 0.6

ny = 400
y = (np.arange(ny) + 0.5) * h / ny                # fibres de béton, y mesuré depuis la base
dy = h / ny
yc = h / 2


def sigma_concrete(eps):
    """Loi du béton de l'article : Desayi en compression (nul après eps_cu), raidissement en traction."""
    eps = np.asarray(eps, dtype=float)
    s = np.zeros_like(eps)
    ecr = fcr / E_c
    t = eps > 0
    lin = t & (eps <= ecr)
    s[lin] = E_c * eps[lin]
    soft = t & (eps > ecr)
    s[soft] = fcr / (1.0 + np.sqrt(500.0 * eps[soft]))
    eu = 1.8 * fc / E_c
    c = (eps < 0) & (-eps <= eps_cu)
    x = -eps[c]
    s[c] = -2.0 * fc * x / (eu * (1.0 + (x / eu) ** 2))
    return s


def sigma_steel(eps):
    return np.clip(E_s * eps, -fy, fy)


def forces(eps0, phi):
    """Effort normal N et moment M (autour du centre de gravité) pour la déformation eps0 et la courbure phi."""
    e = eps0 + phi * (y - yc)                       # phi > 0 : fibre inférieure tendue
    sc = sigma_concrete(e)
    N = np.sum(sc) * b * dy
    M = np.sum(sc * (y - yc)) * b * dy
    for d, A in rebars:
        es = eps0 + phi * (d - yc)
        ss = sigma_steel(es)
        sc_lost = sigma_concrete(np.array([es]))[0]  # le béton déplacé par la barre
        N += (ss - sc_lost) * A
        M += (ss - sc_lost) * A * (d - yc)
    return N, M


def moment_curvature(phis):
    Ms = []
    e0 = 0.0
    for phi in phis:
        lo, hi = -0.02, 0.02
        for _ in range(60):                          # dichotomie sur eps0 : N(eps0) croissant
            mid = 0.5 * (lo + hi)
            if forces(mid, phi)[0] > 0:
                hi = mid
            else:
                lo = mid
        e0 = 0.5 * (lo + hi)
        Ms.append(forces(e0, phi)[1])
    return np.array(Ms)


def load_deflection():
    phis = np.linspace(1e-6, 0.06, 600)
    Ms = moment_curvature(phis)
    imax = int(np.argmax(Ms))                        # on reste avant le pic du moment
    phis, Ms = phis[:imax + 1], Ms[:imax + 1]
    # courbure le long de la demi-poutre (symétrie), moment de flexion 4 points M(x) = P/2 * min(x, a)
    xs = np.linspace(0.0, L / 2, 201)
    shape = np.minimum(xs, a) / 2.0                  # M(x) / P
    m_unit = xs / 2.0                                # moment dû à une charge unitaire à mi-portée
    Ps = np.linspace(0.0, Ms[-1] / (a / 2.0), 300)
    out = []
    for P in Ps:
        kappa = np.interp(P * shape, Ms, phis)       # courbure en chaque x
        delta = 2.0 * np.trapezoid(kappa * m_unit, xs)   # théorème de la charge unitaire
        out.append((delta, P))
    out = np.array(out)
    return out[:, 0] * 1000.0, out[:, 1] / 1000.0, phis, Ms


if __name__ == "__main__":
    d, P, phis, Ms = load_deflection()
    # rigidité initiale, moment de fissuration et moment/charge maximaux
    k0 = (P[5] / d[5])
    icr = np.argmax(np.abs(np.gradient(P, d)) < 0.6 * np.gradient(P, d)[3])
    print(f"raideur initiale (section)      : {k0:.1f} kN/mm")
    print(f"charge à la rupture de pente   : {P[icr]:.1f} kN a {d[icr]:.2f} mm")
    print(f"moment max de la section       : {Ms.max() / 1000:.1f} kN.m  ->  P = {Ms.max() / (a / 2) / 1000:.1f} kN")
    np.savetxt("beton_section.csv", np.column_stack([d, P]), delimiter=",", header="fleche_mm,P_kN", comments="")
    try:
        import matplotlib
        matplotlib.use("Agg")
        import matplotlib.pyplot as plt
        fig, ax = plt.subplots(figsize=(6.4, 4.0))
        ax.plot(d, P, "k--", lw=1.4, label="analyse de section (RDM)")
        if len(sys.argv) > 1:
            r = np.genfromtxt(sys.argv[1], delimiter=",", names=True)
            ax.plot(r["fleche_mm"], r["P_kN"], color="tab:blue", lw=1.4, label="MPM")
        ax.set_xlabel("flèche à mi-portée (mm)")
        ax.set_ylabel("charge P (kN)")
        ax.set_xlim(0, 26)
        ax.set_ylim(0, 160)
        ax.grid(alpha=0.3)
        ax.legend()
        fig.tight_layout()
        fig.savefig("tuto_beton_courbe.png", dpi=170)
        print("figure : tuto_beton_courbe.png")
    except ImportError:
        print("matplotlib absent : pas de figure")
\end{lstlisting}

\section*{Référence}

Cremona C., Houde M.-J., \emph{Modélisation déterministe et probabiliste du comportement mécanique
simplifié des corps d'épreuve}, RFGC, projet RGC\&U \og benchmark des poutres de la Rance \fg{}.

\end{document}
